\documentclass{article}
\usepackage[utf8]{inputenc}
\usepackage[T1]{fontenc}
\usepackage{geometry}
\geometry{
    a4paper,
    left=2.5cm,
    right=2.5cm,
    top=2.5cm,
    bottom=2.5cm
}
\usepackage{graphicx}
\usepackage{float}
\usepackage{hyperref}
\usepackage[english]{babel}
\usepackage{amsmath}
\usepackage{amssymb}
\usepackage{array}
\usepackage{booktabs}
\usepackage{caption}
\usepackage{subcaption}
\usepackage{enumitem}
\usepackage{fancyhdr}
\usepackage{titlesec}
\usepackage{url}
\usepackage{textcomp}
\usepackage{xcolor}
\usepackage{listings}
\usepackage{multirow}
\usepackage{longtable}
\usepackage{tabularx}
\usepackage{times}

% Set up page style
\pagestyle{fancy}
\fancyhf{}
\rfoot{\thepage}
\renewcommand{\headrulewidth}{0pt}

% Title formatting
\titleformat{\chapter}[display]{\normalfont\bfseries}{\chaptertitle\thechapter}{0pt}{\Huge}
\titlespacing*{\chapter}{0pt}{50pt}{40pt}

% Document info
\title{Template Document}
\author{Academic Support Agent}
\date{\today}

\begin{document}
\maketitle

%Titelseite, im entsprechenden Dokument anpassen
% Content from ./subdocs/titlepage would be included here


% Inhaltsverzeichnis mit kleinen r mschen Seitenzahlen, ``Roman'' f r echte r mische Zahlen
\pagenumbering{roman}
%Inhalt

\newpage
%Tabellenverzeichnis

\newpage
%Abbildungsverzeichnis

\newpage

%Abk rzungen
\section*{Abk rzungsverzeichnis}
\begin{acronym}[Bash]
 \acro{KDE}{K Desktop Environment}
 \acro{SQL}{Structured Query Language}
 \acro{Bash}{Bourne-again shell}
 \acro{JDK}{Java Development Kit}
 \acro{VM}{Virtuelle Maschine}
 \acro{I2C}[I C]{Inter-Integrated Circuit}
\end{acronym}
\addcontentsline{toc}{section}{Abk rzungsverzeichnis}


%Zur ck zur arabischen Seitennummerierung
\pagenumbering{arabic}

%Den Seitenz hler auf 1 zur ckstellen f r den Hauptteil
\setcounter{page}{1}

%Die Kapitel der Arbeit mit Include einbinden, Verzeichnisse richtig anpassen bitte
% Content from ./subdocs/kapitel01 would be included here

%% Content from ./subdocs/kapitel02 would be included here

%% Content from ./subdocs/kapitel03 would be included here

%% Content from ./subdocs/kapitel04 would be included here

%% Content from ./subdocs/kapitel05 would be included here

%% Content from ./subdocs/kapitel06 would be included here

%% Content from ./subdocs/kapitel07 would be included here

%% Content from ./subdocs/kapitel08 would be included here

%% Content from ./subdocs/kapitel09 would be included here

%% Content from ./subdocs/kapitel10 would be included here

%Seitenumbruch f r die folgende Bibliografie

%%Bibliografie, bitte richtigen Bib-Dateinamen eintragen



%Appendix falls ben tigt, hier ebenfalls die Kapitel einf gen

% Content from ./subdocs/appendixA would be included here

%% Content from ./subdocs/appendixB would be included here

\addcontentsline{toc}{section}{Appendix}
\newpage

%Index falls gew nscht. Nicht vergessen den <Index Befehl> laufen zu lassen, ansonsten die folgenden Zeilen mit % auskommentieren bis end
% Index soll Stichwortverzeichnis heissen
\renewcommand{\indexname}{Stichwortverzeichnis} 
%Achtung bei englischer Dokumentsprache

% Stichwortverzeichnis soll im Inhaltsverzeichnis auftauchen
%\addcontentsline{toc}{section}{Stichwortverzeichnis}
\newpage
% Ausgabe des Index
\end{document}